% concatenated from main.tex
\documentclass[11pt]{article}

% -------------
% Packages
% -------------
\usepackage[margin=1in]{geometry}
\usepackage{amsmath,amssymb,amsthm,mathtools,bm,bbm}
\usepackage[authoryear]{natbib}
\usepackage{hyperref}
\hypersetup{
    colorlinks=true,
    linkcolor=blue,
    citecolor=blue,
    urlcolor=blue
}


\theoremstyle{plain}
\newtheorem{theorem}{Theorem}
\newtheorem{lemma}[theorem]{Lemma}
\newtheorem{corollary}[theorem]{Corollary}
\newtheorem{assumption}[theorem]{Assumption}
\newtheorem{proposition}[theorem]{Proposition}

\theoremstyle{remark}
\newtheorem{remark}[theorem]{Remark}
\DeclareMathOperator{\Var}{Var}
\DeclareMathOperator{\Tr}{Tr}
\DeclareMathOperator{\spec}{spec}
\DeclareMathOperator{\dist}{dist}
\newcommand{\E}{\mathbb{E}}
\newcommand{\PP}{\mathbb{P}}
\newcommand{\Unif}{\mathrm{Unif}}
\newcommand{\op}{\mathrm{op}}
\newcommand{\HS}{\mathrm{HS}}
\newcommand{\II}{[0,1]^2}


\title{\textbf{Distributional Limits for Eigenvalues of Graphon Kernel Matrices}}

\author{
  Behzad Aalipur\thanks{
    Department of Mathematical Sciences, University of Cincinnati.
    \texttt{aalipubd@mail.uc.edu}
  }
}
\date{}


\begin{document}

\maketitle

\begin{abstract}
We study the fluctuation behavior of individual eigenvalues of kernel matrices arising from dense graphon-based random graphs. Under minimal integrability and boundedness assumptions on the graphon, we establish distributional limits for simple, well-separated eigenvalues of the associated integral operator. A sharp probabilistic dichotomy emerges: in the non-degenerate regime, the properly normalized empirical eigenvalue satisfies a central limit theorem with an explicit variance, whereas in the degenerate regime the leading stochastic term vanishes and the centered eigenvalue converges to a weighted chi-square law determined by the operator spectrum.

The analysis requires no smoothness or Lipschitz conditions on the kernel. Prior work under comparable assumptions established only operator convergence and eigenspace consistency; the present results characterize the full distributional behavior of individual eigenvalues, extending fluctuation theory beyond the reach of classical operator-level arguments. The proofs combine second-order perturbation expansions, concentration bounds for kernel matrices, and Hoeffding decompositions for symmetric statistics, revealing that at the $\sqrt{n}$ scale the dominant randomness arises from latent-position sampling rather than Bernoulli edge noise.
\end{abstract}

\textbf{Keywords:} Graphon models; eigenvalue distributions; kernel matrices; U-statistics; spectral perturbation.

\textbf{MSC2020:} Primary 60B20; Secondary 60F05, 62G20, 05C80.


\section{Introduction}

Spectral methods play a central role in the analysis of random graphs and network data 
\cite{lei2015consistency,rohe2011spectral}. Their success in community detection 
\cite{abbe2017community,lei2015consistency,zhao2012consistency}, clustering 
\cite{von2007tutorial}, and goodness-of-fit procedures \cite{bickel2011method,gao2015rate} 
has motivated a substantial body of work on the asymptotic behavior of eigenvalues and 
eigenvectors in latent-variable models. Classical results establish almost-sure or $L^2$ 
convergence of empirical kernel operators \cite{koltchinskii2000random,karoui2008spectrum,braun2006accurate}, 
together with consistency of eigenspaces and certain linear spectral statistics. 
However, the distributional behavior of individual eigenvalues, a question of central 
importance for uncertainty quantification remains comparatively underdeveloped.

This paper develops a distributional theory for simple eigenvalues of kernel matrices generated by dense graphon-based random graphs. Working within the graphon framework of Lovász and Szegedy \cite{lovasz2006limits,lovasz2012large,borgs2008convergent,borgs2012convergent}, we show that each empirical eigenvalue exhibits a sharp probabilistic dichotomy in its fluctuations. Depending on the degeneracy 
structure of the population eigenfunction, the leading fluctuations are either asymptotically 
Gaussian or converge to an explicit weighted chi-square law. A version of this dichotomy was first obtained by \cite{chatterjee2025fluctuation} under Lipschitz regularity of the graphon. Our contribution is to extend this phenomenon to the minimal integrability setting 
$W\in L^2([0,1]^2)\cap L^\infty([0,1]^2)$, thereby covering discontinuous, non-smooth, and 
piecewise-constant models such as stochastic block models with unequal block sizes. The results 
demonstrate that smoothness assumptions, while technically convenient, are not intrinsic to the 
limiting behavior.

From a probabilistic perspective, the fluctuation behavior is driven by the interaction between 
two sources of randomness: latent-position sampling and Bernoulli edge noise. A central finding 
of this work is that, at the $\sqrt{n}$ scale, the dominant contribution arises from the latent 
variables, while edge noise is asymptotically negligible. This separation, implicit in earlier 
work on spectral convergence \cite{koltchinskii2000random} and latent-position 
models \cite{lei2015consistency,abbe2017community,athreya2018statistical}, becomes explicit in 
our analysis and clarifies the sources of uncertainty in spectral procedures. In particular, the 
zero diagonal of the kernel matrix-an intrinsic structural feature of graphon models-plays a 
decisive role in the degenerate regime, where the linear Hoeffding projection vanishes and the 
limiting distribution becomes a weighted chi-square series.

Our proof combines three probabilistic ingredients. First, a second-order Rayleigh-Schrödinger expansion for simple eigenvalues (\cite{kato1966perturbation}, Lemma ~\ref{lem:expansion}) reduces the problem to the analysis 
of a centered Rayleigh quotient. Second, concentration of kernel matrices under boundedness 
assumptions yields operator-norm control of the empirical perturbation. Third, a Hoeffding 
decomposition for the resulting symmetric statistic \cite{serfling2009approximation} isolates the 
linear and degenerate components, revealing the mechanism behind the Gaussian versus weighted 
chi-square regimes. This synthesis demonstrates that the fluctuation dichotomy persists well beyond the smooth settings considered in \cite{chatterjee2025fluctuation}, and that the limiting distributions 
are determined by spectral structure and degeneracy rather than regularity of the kernel.

The remainder of the paper is organized as follows. Section~\ref{sec:prelim} introduces the 
graphon model, the kernel matrix, and the associated integral operator. Section~\ref{sec:main} 
states the main fluctuation theorem and outlines the perturbation-theoretic approach. 
Section~\ref{sec:examples} presents non-Lipschitz examples illustrating the scope of the theory. 
Section~\ref{sec:discussion} discusses implications for spectral methods and directions for future 
work. Technical proofs are given in Section~\ref{sec:proof_strategy}.


\section{Preliminaries}
\label{sec:prelim}

A graphon is a measurable symmetric function
\begin{equation}
W : [0,1]^2 \to [0,1], \quad W(x,y)=W(y,x) \text{ for almost every } (x,y)\in[0,1]^2.
\label{eq:graphon_def}
\end{equation}
Graphons provide a canonical framework for modeling dense inhomogeneous random graphs and arise as limits of sequences of finite graphs
\citep{lovasz2006limits,lovasz2012large,borgs2008convergent,borgs2012convergent}.
To generate a random graph on $n$ vertices, we first sample latent variables
\begin{equation}
U_1,\dots,U_n \stackrel{\mathrm{i.i.d.}}{\sim} \mathrm{Unif}[0,1].
\label{eq:latent_vars}
\end{equation}
Conditional on $(U_1,\dots,U_n)$, edges are generated independently according to
\begin{equation}
A_{ij} \mid U_1,\dots,U_n \;\sim\; \mathrm{Bernoulli}\bigl(W(U_i,U_j)\bigr), \quad i<j,
\label{eq:edge_generation}
\end{equation}
with $A_{ji}=A_{ij}$ and $A_{ii}=0$.
This defines a simple undirected random graph whose conditional expectation satisfies
\begin{equation}
\mathbb{E}\bigl[\mathbf{A}_n \mid U_1,\dots,U_n\bigr] = \mathbf{K}_n,
\label{eq:conditional_mean}
\end{equation}
where $\mathbf{K}_n$ denotes the kernel matrix defined below
\citep{bickel2009nonparametric,wolfe2013nonparametric,chatterjee2015matrix}.

The kernel matrix is given by
\begin{equation}
(\mathbf{K}_n)_{ij} :=
\begin{cases}
W(U_i,U_j), & i\neq j,\\
0, & i=j.
\end{cases}
\label{eq:kernel_matrix}
\end{equation}
The zero diagonal reflects the absence of self-loops and is intrinsic to graphon-based network models.
This structural restriction plays a nontrivial role in the fluctuation behavior, particularly in degenerate regimes.
As will be seen later, it has important consequences for the asymptotic behavior of spectral statistics.


The graphon $W$ induces a compact self-adjoint integral operator
\begin{equation}
T_W : L^2([0,1]) \to L^2([0,1]), \quad (T_W f)(x) := \int_0^1 W(x,y)f(y)\,dy.
\label{eq:integral_operator}
\end{equation}
Since $W\in L^2([0,1]^2)$, the operator $T_W$ is Hilbert-Schmidt and therefore compact.
Let $(\lambda_k,\varphi_k)_{k\ge1}$ denote its eigenvalues and eigenfunctions, ordered so that $|\lambda_1|\ge|\lambda_2|\ge\cdots\ge0$ and normalized according to
\begin{equation}
\mathbb{E}[\varphi_k(U)^2]=1, \quad U\sim\mathrm{Unif}[0,1].
\label{eq:eigenfunction_normalization}
\end{equation}
For each fixed $k$, the empirical eigenvalues of the kernel matrix satisfy
\begin{equation}
\lambda_k(\mathbf{K}_n)/(n-1)\xrightarrow{\mathbb{P}}\lambda_k
\quad\text{as }n\to\infty,
\label{eq:eigenvalue_convergence}
\end{equation}
see for example \citet{koltchinskii2000random,zhu2024central}. 

The corresponding population eigenfunction $\varphi_r$ is discretized as
\begin{equation}
\mathbf{v}_r := (\varphi_r(U_1),\dots,\varphi_r(U_n))^\top\in\mathbb{R}^n.
\label{eq:unnormalized_eigenvector}
\end{equation}
Define the normalization factor
\begin{equation}
s_{r,n} := \sum_{i=1}^n\varphi_r(U_i)^2 = n\bigl(1+V_{r,n}\bigr),
\label{eq:normalization_factor}
\end{equation}
where
\begin{equation}
V_{r,n} := \frac{1}{n}\sum_{i=1}^n\bigl(\varphi_r(U_i)^2-1\bigr).
\label{eq:centered_eigenfunction_variance}
\end{equation}
The normalized discretization vector is
\begin{equation}
\mathbf{u}_r := s_{r,n}^{-1/2}(\varphi_r(U_1),\dots,\varphi_r(U_n))^\top,
\label{eq:normalized_eigenvector}
\end{equation}
which satisfies $\|\mathbf{u}_r\|_2=1$.

Let $A$ be a self-adjoint operator with a simple eigenvalue $\lambda_r$ and normalized eigenvector $\psi_r$, and let $E$ be a self-adjoint perturbation. The second-order Rayleigh-Schr\"odinger expansion yields
\[
\lambda_r(A+E)-\lambda_r(A)
=
\langle\psi_r,E\psi_r\rangle
+
\sum_{k\neq r}\frac{|\langle\psi_k,E\psi_r\rangle|^2}{\lambda_r-\lambda_k}
+
R(E),
\]
with remainder bounded by $|R(E)|\le 2\|E\|_{\mathrm{op}}^2/\gamma_r$.
In the present setting, the leading term corresponds to the discrete Rayleigh quotient
\begin{equation}
\mathbf{u}_r^\top\mathbf{K}_n\mathbf{u}_r
=
\frac{1}{s_{r,n}}
\sum_{i\neq j}\varphi_r(U_i)W(U_i,U_j)\varphi_r(U_j).
\label{eq:rayleigh_quotient}
\end{equation}
This quantity is a second-order U-statistic with symmetric kernel
\begin{equation}
h_r(x,y) := \varphi_r(x)W(x,y)\varphi_r(y).
\label{eq:ustat_kernel}
\end{equation}
The Hoeffding decomposition gives
\begin{equation}
\frac{1}{n(n-1)}\sum_{i\neq j}h_r(U_i,U_j)
=
\theta
+
\frac{2}{n}\sum_{i=1}^n h_1(U_i)
+
U_n^{(2)},
\label{eq:hoeffding_decomposition}
\end{equation}
where
\(
\theta := \mathbb{E}[h_r(U,U')]\), \( h_1(x) := \mathbb{E}[h_r(x,U)]-\theta\), and
\(U_n^{(2)}\) is a degenerate second-order U-statistic.



We defer discussion of how the linear and degenerate Hoeffding components determine the Gaussian versus weighted chi-square regimes to Section~\ref{sec:proof_strategy}.


\section{Main Results}
\label{sec:main}

Recall from Section~\ref{sec:prelim} that $\mathbf{K}_n$ denotes the kernel matrix with zero diagonal,
and $\mathbf{A}_n$ is the adjacency matrix with conditional mean $\mathbb{E}[\mathbf{A}_n \mid U_1,\dots,U_n] = \mathbf{K}_n$.

To isolate the leading-order fluctuations of a target eigenvalue $\lambda_r$, we work under a small 
set of regularity conditions on the kernel $W$ and its spectrum. These assumptions are very mild and 
designed to ensure that the integral operator $T_W$ has a simple, well-separated eigenvalue with an eigenfunction having finite fourth moment that are natural in the graphon setting and hold for a broad class 
of non-Lipschitz kernels.

\begin{assumption}\label{ass:ustat}
Fix an index $r\ge1$. Let $W\in L^2([0,1]^2)\cap L^\infty([0,1]^2)$ be symmetric, 
and let $T_W:L^2([0,1])\to L^2([0,1])$ denote the associated integral operator
\[
(T_W f)(x) := \int_0^1 W(x,y) f(y)\,dy.
\]
Let $(\lambda_k, \varphi_k)_{k\ge1}$ denote an $L^2$-orthonormal eigendecomposition of $T_W$, 
ordered by decreasing $|\lambda_k|$. Assume:

\begin{enumerate}
\item (\textbf{Spectral gap at $r$}) $\lambda_r\neq 0$ is simple and
\[
\gamma_r:=\min_{k\neq r}|\lambda_r-\lambda_k|>0.
\]
\item (\textbf{Finite fourth moment}) $\E[\varphi_r(U)^4]<\infty$.

\end{enumerate}
\end{assumption}

\begin{remark}
Since $T_W$ is Hilbert-Schmidt, each eigenfunction $\varphi_k$ lies in $L^2([0,1])$ automatically;
we only impose additional moment condition stated above when needed for distributional limit theorems.
\end{remark}


\begin{theorem}[Unified eigenvalue fluctuation dichotomy]
\label{thm:ustat}
Under Assumption~\ref{ass:ustat}, let $\lambda_r$ denote the $r$th eigenvalue of the 
population integral operator $T_W$. Then:

\begin{enumerate}
\item \textbf{Non-degenerate case.}
If $\sigma_r^2 := \Var(\varphi_r(U)^2) > 0$, then
\[
\sqrt{n}\left(\frac{\lambda_r(\mathbf{K}_n)}{n-1}-\lambda_r\right)
\xrightarrow{d}
\mathcal{N}(0,\lambda_r^2\sigma_r^2).
\]

\item \textbf{Degenerate case.}
If $\sigma_r^2 := \Var(\varphi_r(U)^2) = 0$, which holds if and only if 
$\varphi_r^2 \equiv 1$ almost everywhere, then
\[
\lambda_r(\mathbf{K}_n) - (n-1)\lambda_r - C_r
\xrightarrow{d}
\sum_{k\neq r}\frac{\lambda_r\lambda_k}{\lambda_r-\lambda_k}(Z_k^2-1),
\qquad
C_r:=\sum_{k\neq r}\frac{\lambda_k^2}{\lambda_r-\lambda_k}.
\]
where $(Z_k)_{k\ge1}$ are i.i.d.\ standard normal random variables, and the series 
converges in $L^2$.
\end{enumerate}
\end{theorem}

\begin{remark}[Origin of the degenerate coefficients]
The coefficient $\frac{\lambda_r \lambda_k}{\lambda_r - \lambda_k}$ is the algebraic sum of 
two distinct sources of fluctuation.
Note that 
\[
\lambda_k(Z_k^2-1)
+
\frac{\lambda_k^2}{\lambda_r-\lambda_k}(Z_k^2-1)
=
\frac{\lambda_r\lambda_k}{\lambda_r-\lambda_k}(Z_k^2-1),
\]
Both terms are $O_p(1)$ in the degenerate regime.

\begin{enumerate}
    \item The \textbf{linear term} contributes $\lambda_k(Z_k^2-1)$.
    This fluctuation arises specifically because $\mathbf{K}_n$ has a zero diagonal; even for the case
    of $\mathbb{E}[h_{r,2}(U_1,U_2)\mid U_1]=0$, the 
    degenerate U-statistic $\frac{1}{n}\sum_{i\neq j}h_{r,2}(U_i,U_j)$ does not vanish.
    \item The \textbf{resolvent term} contributes 
    $\frac{\lambda_k^2}{\lambda_r - \lambda_k}(Z_k^2-1)$ due to the spectral perturbation 
    from orthogonal modes interacting through the Rayleigh-Schrödinger expansion.
\end{enumerate}

\end{remark}


\subsection{Proof strategy}
\label{sec:proof_strategy}

We outline the main steps in the proof of Theorem~\ref{thm:ustat}.
The analysis is carried out at the matrix level using second-order
perturbation theory for simple eigenvalues, combined with Hoeffding
decompositions of U-statistics.

We work with the normalized kernel matrix
\[
\widetilde{\mathbf K}_n := \frac{1}{n-1}\mathbf K_n,
\qquad
\widetilde{\mathbf\Delta}_n
:= \widetilde{\mathbf K}_n - \mathbb E[\widetilde{\mathbf K}_n].
\]
By Lemma~\ref{lem:approx_eigvec}, the $r$th eigenvector of
$\mathbb E[\widetilde{\mathbf K}_n]$ is well approximated by the
normalized discretization $\mathbf u_r$, and the corresponding eigengap
is bounded below by $\gamma_r$. Applying the Rayleigh-Schrödinger
expansion for simple eigenvalues yields
\begin{equation}
\lambda_r(\widetilde{\mathbf K}_n)-\lambda_r
=
\mathbf u_r^\top \widetilde{\mathbf\Delta}_n \mathbf u_r
+
\text{\emph{(resolvent correction)}}
+
R_n,
\label{eq:rs_outline}
\end{equation}
where the remainder $R_n$ is quadratic in
$\widetilde{\mathbf\Delta}_n$ and satisfies $R_n=O_p(n^{-1})$ under the
boundedness of $W$, and is therefore negligible at the $\sqrt n$ scale.

The leading term in \eqref{eq:rs_outline} is a centered Rayleigh quotient,
which can be written as a second-order U-statistic with kernel
$h_r(x,y)=\varphi_r(x)W(x,y)\varphi_r(y)$. Its Hoeffding decomposition
isolates a linear component proportional to
$V_{r,n}=\frac1n\sum_{i=1}^n(\varphi_r(U_i)^2-1)$.
If $\sigma_r^2=\Var(\varphi_r(U)^2)>0$, this linear term dominates and
the classical central limit theorem yields Gaussian fluctuations.

The resolvent correction in \eqref{eq:rs_outline} corresponds to the
second-order term in the Rayleigh-Schrödinger expansion,
\[
\sum_{k\neq r}
\frac{
(\mathbf u_k^\top \widetilde{\mathbf\Delta}_n \mathbf u_r)^2
}{
\lambda_r-\lambda_k
},
\]
which captures the effect of random projections of the target eigenvector
onto orthogonal eigenspaces, weighted by the inverse spectral gaps.
In the degenerate regime $\varphi_r^2\equiv1$, the linear Hoeffding term
vanishes identically, and the leading fluctuations arise from the
degenerate U-statistic component of the Rayleigh quotient together with
this resolvent correction.

Both contributions can be expressed in terms of the empirical
cross-projections
\[
T_{k,n}=\frac{1}{\sqrt n}\sum_{i=1}^n
\varphi_r(U_i)\varphi_k(U_i),
\qquad k\neq r.
\]
Combining these terms yields weighted quadratic forms in
$T_{k,n}^2-1$, and leads to the non-Gaussian weighted chi-square limit in
the degenerate case. Full technical details are deferred to
Section~\ref{sec:tech}.



\subsection{Implications for the adjacency matrix}
When analyzing the adjacency matrix $\mathbf{A}_n$ with conditional mean $\mathbb{E}[\mathbf{A}_n\mid U_1,\dots,U_n]=\mathbf{K}_n$, 
the eigenvalue fluctuations at the $\sqrt{n}$ scale are governed by the latent kernel $\mathbf{K}_n$, not by the Bernoulli edge noise.
\begin{corollary}[Adjacency eigenvalue fluctuations in the non-degenerate regime]
\label{cor:edge_noise}
Let $\mathbf{A}_n$ be the adjacency matrix with $\mathbb{E}[\mathbf{A}_n\mid U_1,\dots,U_n]=\mathbf{K}_n$.
Assume Assumption~\ref{ass:ustat} and suppose $\sigma_r^2:=\Var(\varphi_r(U)^2)>0$.
Then
\[
\sqrt{n}\left(
\lambda_r\!\left(\frac{\mathbf{A}_n}{n-1}\right)
-
\lambda_r\!\left(\frac{\mathbf{K}_n}{n-1}\right)
\right)
\overset{p}{\longrightarrow} 0,
\]
and
\[
\sqrt{n}\left(
\lambda_r\!\left(\frac{\mathbf{A}_n}{n-1}\right)
-
\lambda_r
\right)
\xrightarrow{d}
\mathcal{N}\!\left(0,\lambda_r^2\sigma_r^2\right).
\]
\end{corollary}

\begin{proof}
The first display follows from the eigenvalue perturbation bound for inhomogeneous Bernoulli graphs; see Theorem~2 in \citet{aalipur2026bootstrap}.
The second display follows by Slutsky's theorem, combining the first display with the non-degenerate conclusion of Theorem~\ref{thm:ustat} for $\lambda_r(\mathbf{K}_n/(n-1))$.
\end{proof}


\begin{remark}
If $\varphi_r^2\equiv 1$ almost everywhere, then the limit in Theorem~\ref{thm:ustat} is non-Gaussian, and the same non-Gaussian limit carries over from $\mathbf{K}_n$ to $\mathbf{A}_n$ under the same edge-noise perturbation bound.
\end{remark}


\section{Examples}
\label{sec:examples}

This section presents several examples of graphons that satisfy Assumption~\ref{ass:ustat} and therefore fall within the scope of Theorem~\ref{thm:ustat}, while lying outside the Lipschitz framework considered in earlier work. These examples illustrate that $\sqrt{n}$-scale eigenvalue fluctuations depend primarily on spectral and integrability properties rather than smoothness, and they clarify the distinction between the non-degenerate and degenerate regimes.

\subsection{Brownian square-root kernel}

Define $W(x,y):=\min\{x,y\}+\sqrt{xy}$ on $[0,1]^2$, and write $W=K_0+K_1$ with
$K_0(x,y):=\min\{x,y\}$ (Brownian covariance kernel, $\lambda_k(K_0)\asymp k^{-2}$)
and $K_1(x,y):=\sqrt{xy}$ (rank one). Hence $T_W$ is a finite-rank perturbation of
$T_{K_0}$ and $\lambda_k(W)\asymp k^{-2}$.

The kernel is not Lipschitz: $\partial_x\sqrt{xy}=\sqrt{y}/(2\sqrt{x})$ diverges as
$x\to0$, and the same local obstruction implies it is not piecewise Lipschitz under
any finite partition.

Nevertheless, $W\in L^2([0,1]^2)\cap L^\infty([0,1]^2)$ with $\|W\|_\infty\le2$.
Under Assumption~\ref{ass:ustat}, Theorem~\ref{thm:ustat} applies and yields
$\sqrt{n}(\lambda_r(\mathbf{K}_n)/(n-1)-\lambda_r)\xrightarrow{d}\mathcal{N}(0,\lambda_r^2\sigma_r^2)$
whenever $\sigma_r^2>0$, showing coverage beyond Lipschitz kernels.



\subsection{Stochastic block model with unequal block sizes}

Consider a two-block stochastic block model with block proportions $\pi_1=1/3$ and $\pi_2=2/3$, within-block edge probability $p$, and cross-block probability $q<p$. The associated graphon is
\[
W(x,y)=
\begin{cases}
p, & (x,y)\in[0,1/3]^2\cup[1/3,1]^2,\\
q, & \text{otherwise}.
\end{cases}
\]
This kernel is piecewise constant and discontinuous. The second eigenfunction $\varphi_2$ encodes the block structure. Due to unequal block sizes, $\varphi_2$ is non-constant and attains larger absolute values on the larger block. Consequently,
$\mathbb{E}[\varphi_2(U)^2]=1$ and $\Var(\varphi_2(U)^2)>0$, placing this model in the non-degenerate regime.

The kernel satisfies Assumption~\ref{ass:ustat}. It belongs to $L^2([0,1]^2)$ and $L^\infty([0,1]^2)$, its spectrum is finite dimensional and determined by the block structure, and its eigenfunctions are bounded and piecewise constant. Theorem~\ref{thm:ustat} therefore yields
\[
\sqrt{n}\left(\frac{\lambda_2(\mathbf{K}_n)}{n-1}-\lambda_2\right)
\xrightarrow{d}
\mathcal{N}\bigl(0,\lambda_2^2\sigma_2^2\bigr),
\]
where $\sigma_2^2=\Var(\varphi_2(U)^2)>0$. This example illustrates that the theory applies to practically relevant network models that fall outside Lipschitz-based analyses.

\subsection{Symmetric stochastic block model}

For a symmetric two-block stochastic block model with equal block sizes $\pi_1=\pi_2=1/2$, within-block probability $p$, and cross-block probability $q<p$, the graphon is
\[
W(x,y)=
\begin{cases}
p, & (x,y)\in[0,1/2]^2\cup[1/2,1]^2,\\
q, & \text{otherwise}.
\end{cases}
\]
By symmetry, the second eigenfunction satisfies $\varphi_2(x)\in\{-1,1\}$ almost everywhere, so that $\varphi_2^2\equiv1$ and $\Var(\varphi_2(U)^2)=0$. This model therefore lies in the degenerate regime of Theorem~\ref{thm:ustat}. The eigenvalue fluctuation is non-Gaussian and satisfies
\[
\lambda_2(\mathbf{K}_n)-(n-1)\lambda_2
\xrightarrow{d}
\sum_{k\neq2}\frac{\lambda_2\lambda_k}{\lambda_2-\lambda_k}(Z_k^2-1),
\]
with convergence in $L^2$ under the Hilbert-Schmidt condition. Despite its structural similarity to the unequal-block model, symmetry induces a fundamentally different limiting distribution, illustrating the dichotomy underlying Theorem~\ref{thm:ustat}.

\subsection{\texorpdfstring{$\alpha$-H\"older kernels}{alpha-Holder kernels}}

Consider the power-law kernel $W(x,y)=(xy)^\alpha$ for $\alpha\in(0,1)$. The kernel satisfies Assumption~\ref{ass:ustat}. It belongs to $L^2([0,1]^2)$ since
\[
\int_0^1\int_0^1(xy)^{2\alpha}\,dx\,dy
=
\left(\int_0^1 x^{2\alpha}\,dx\right)^2
=
\frac{1}{(2\alpha+1)^2}<\infty,
\]
and it belongs to $L^\infty([0,1]^2)$ with $\|W\|_\infty=1$. The associated operator is compact and positive definite, and its eigenfunctions are continuous and bounded.

The kernel is not Lipschitz. Its partial derivatives satisfy
\[
\frac{\partial W}{\partial x}(x,y)=\alpha y^\alpha x^{\alpha-1},
\]
which diverges as $x\to0$ for $\alpha<1$. Consequently, a Lipschitz bound of the form
\[
|x^\alpha-x'^\alpha|\le L|x-x'|
\]
cannot hold uniformly on $[0,1]$. The leading eigenfunction $\varphi_1$ is strictly positive and non-constant, so that $\Var(\varphi_1(U)^2)>0$. Theorem~\ref{thm:ustat} therefore applies in the non-degenerate regime, yielding a Gaussian central limit theorem for $\lambda_1(\mathbf{K}_n)/(n-1)$.


\section{Discussion}
\label{sec:discussion}

This work develops a distributional theory for individual eigenvalues of dense graphon-based random graphs under minimal $L^2([0,1]^2) \cap L^\infty ([0,1]^2)$ assumptions. Our main result extends the eigenvalue fluctuation dichotomy established by \citet{chatterjee2025fluctuation} beyond Lipschitz kernels and shows that smoothness, while technically convenient in prior analyses, is not intrinsic to the limiting behavior described in Theorem~\ref{thm:ustat}. The key probabilistic mechanism is instead the degeneracy structure induced by the population eigenfunction and the zero-diagonal kernel-matrix normalization.

A central conceptual takeaway is that, at the $\sqrt{n}$ scale, eigenvalue fluctuations are driven primarily by sampling variability in the latent positions rather than by Bernoulli edge noise. In the non-degenerate regime, where $\sigma_r^2=\Var(\varphi_r(U)^2)>0$, this leads to a Gaussian limit with variance $\lambda_r^2\sigma_r^2$, as in Theorem~\ref{thm:ustat}. In the degenerate regime, where $\varphi_r^2\equiv 1$ almost everywhere, the $\sqrt{n}$-scale fluctuations vanish and the centered eigenvalue converges to an explicit weighted chi-square series determined by the spectrum of $T_W$. In stochastic block models, this dichotomy cleanly separates unequal-block settings (typically non-degenerate) from symmetric equal-block settings (degenerate), indicating that the non-Gaussian limit reflects an underlying structural symmetry that is invisible to first-order fluctuations.

From an inferential perspective, the results clarify the sources of uncertainty in spectral procedures for network analysis. In the non-degenerate regime, the asymptotic variance of $\lambda_r(\mathbf K_n)/(n-1)$ is $\lambda_r^2\sigma_r^2/n$, suggesting plug-in estimation strategies for confidence intervals and hypothesis tests for eigenvalue-based methods in latent-position and block models. Moreover, Corollary~\ref{cor:edge_noise} shows that, after normalization, Bernoulli edge noise is asymptotically negligible at the $\sqrt{n}$ scale, so the same distributional limits apply to the observed adjacency matrix $\mathbf A_n$. The degenerate regime offers complementary diagnostic value: detecting $\varphi_r^2\equiv 1$ corresponds to identifying symmetry in the generating structure and can inform model checking and selection.

Several extensions merit further investigation. The approach should extend to joint fluctuations of finitely many simple eigenvalues in the non-degenerate regime, yielding multivariate Gaussian limits with explicit covariance structure. Beyond eigenvalues, combining perturbation expansions with Davis-Kahan-type control may yield fluctuation results for eigenvectors and derived quantities relevant to spectral clustering uncertainty quantification. Sparse and semi-dense regimes, where the graphon or its scaling depends on $n$, require different concentration tools and may exhibit phase transitions in eigenvalue behavior. Additional challenges arise when the target eigenvalue has multiplicity greater than one, in which case perturbation analysis must be carried out at the eigenspace level and the limiting laws are expected to involve weighted quadratic forms over that eigenspace. Finally, it would be natural to adapt the framework to other graph-associated operators, such as normalized Laplacians, though controlling randomness introduced by degree normalization presents further technical difficulties.

Overall, removing smoothness assumptions provides both a technical extension and a conceptual clarification. The persistence of the fluctuation dichotomy under $L^2([0,1]^2) \cap L^\infty ([0,1]^2)$ conditions shows that the limiting distributions are determined by spectral structure and degeneracy rather than kernel regularity. The examples in Section~\ref{sec:examples} highlight that discontinuous block models, non-Lipschitz H\"older kernels, and Brownian-type covariance structures fall within the scope of the theory, substantially enlarging the class of models for which eigenvalue uncertainty quantification is available.



\section{Proof of the main theorem}
\label{sec:tech}
This section contains the proof of Theorem~\ref{thm:ustat}. We follow the strategy outlined in Section~\ref{sec:proof_strategy}: first reduce the eigenvalue fluctuation to a Rayleigh quotient (up to a negligible error) using perturbation theory for simple eigenvalues, and then analyze the resulting quadratic form via a Hoeffding decomposition. The argument splits into two regimes depending on whether the linear Hoeffding projection is nontrivial ($\sigma_r^2>0$) or vanishes identically ($\sigma_r^2=0$).

\begin{proof}[Proof of Theorem~\ref{thm:ustat}]
We treat the non-degenerate and degenerate cases separately, using spectral perturbation theory and U-statistic decompositions.

\textbf{Case 1: Non-degenerate case ($\sigma_r^2 > 0$)} Recall the normalized discretization of eigenfunction and the corresponding scaling as defined in \eqref{eq:normalized_eigenvector} and \eqref{eq:normalization_factor}:,  
\[
\mathbf{u}_r := s_{r,n}^{-1/2}(\varphi_r(U_1),\dots,\varphi_r(U_n))^\top, \qquad
s_{r,n} := \sum_{i=1}^n \varphi_r(U_i)^2.
\]
By Lemma~\ref{lem:approx_eigvec} we have 
\[
\mathbf{K}_n \mathbf{u}_r = (n-1)\lambda_r \mathbf{u}_r + \mathbf{e}_r,\qquad
\|\mathbf{e}_r\|_2 = O_p(\sqrt{n}),\qquad
\|\mathbf{u}_r\|_2=1.
\]
Define the Rayleigh quotient $\eta_n := \mathbf{u}_r^\top \mathbf{K}_n \mathbf{u}_r$ and residual
$\mathbf{r}_n := \mathbf{K}_n\mathbf{u}_r - \eta_n\mathbf{u}_r$. Then
\[
\eta_n
=
\mathbf{u}_r^\top \mathbf{K}_n \mathbf{u}_r
=
(n-1)\lambda_r + \mathbf{u}_r^\top \mathbf{e}_r
=
(n-1)\lambda_r + O_p(\sqrt n),
\]
and
\[
\mathbf{r}_n
=
\mathbf{e}_r - (\mathbf{u}_r^\top\mathbf{e}_r)\mathbf{u}_r,\qquad
\|\mathbf{r}_n\|_2 \le 2\|\mathbf{e}_r\|_2 = O_p(\sqrt n).
\]

By Lemma~\ref{lem:matconc} and the eigengap condition $\gamma_r>0$ for $T_W$, the eigenvalues
of $\mathbf{K}_n/(n-1)$ can be indexed so that
$\lambda_k(\mathbf{K}_n)/(n-1)\xrightarrow{p}\lambda_k$ for each fixed $k$. In particular, with
probability tending to $1$, all eigenvalues of $\mathbf{K}_n$ other than $\lambda_r(\mathbf{K}_n)$
lie at distance at least $c n$ from $(n-1)\lambda_r$ for some $c>0$. On this event, set
\[
\alpha_n := (n-1)\lambda_r - \frac{c n}{2},\qquad
\beta_n := (n-1)\lambda_r + \frac{c n}{2},
\]
so that for all large $n$,
since $\eta_n=(n-1)\lambda_r+o_p(\sqrt{n})$, we have
\[
\alpha_n < \eta_n < \beta_n \text{\,\,and\,\,}
(\alpha_n,\beta_n)\ \text{contains exactly the eigenvalue } \lambda_r(\mathbf{K}_n).
\]
Moreover on this high-probability event, $\eta_n-\alpha_n\asymp n$, $\beta_n-\eta_n\asymp n$. The Kato-Temple inequality for
self-adjoint matrices (see, e.g., Kato~\citet{kato1966perturbation} or \ref{lem:kato-temple}) then yields
\[
\eta_n - \frac{\|\mathbf{r}_n\|_2^2}{\eta_n-\alpha_n}
\;\le\;
\lambda_r(\mathbf{K}_n)
\;\le\;
\eta_n + \frac{\|\mathbf{r}_n\|_2^2}{\beta_n-\eta_n},
\]
so, since $\eta_n-\alpha_n$ and $\beta_n-\eta_n$ are of order $n$ and
$\|\mathbf{r}_n\|_2^2=O_p(n)$,
\[
|\lambda_r(\mathbf{K}_n)-\eta_n|
\le
\frac{\|\mathbf{r}_n\|_2^2}{c'n}
=
O_p(1) = o_p(\sqrt n).
\]
Thus
\[
\lambda_r(\mathbf{K}_n)
=
\eta_n + O_p(1)
=
(n-1)\lambda_r + O_p(\sqrt n).
\] In particular,
\[ \lambda_r(\mathbf{K}_n) = \mathbf{u}_r^\top\mathbf{K}_n\mathbf{u}_r + O_p(\sqrt n), \]

so the eigenvalue can be analyzed via the Rayleigh quotient at $\mathbf{u}_r$.

Using the kernel representation $(\mathbf{K}_n)_{ij} = W(U_i,U_j)\mathbbm{1}\{i \neq j\}$,
\[
\mathbf{u}_r^\top \mathbf{K}_n \mathbf{u}_r
= \frac{1}{s_{r,n}} \sum_{i \neq j} \varphi_r(U_i) W(U_i, U_j) \varphi_r(U_j)
= \frac{1}{s_{r,n}} \sum_{i \neq j} h_r(U_i, U_j),
\]
where $h_r(x,y) := \varphi_r(x)W(x,y)\varphi_r(y)$. For the symmetric kernel $h_r$, the Hoeffding
decomposition (see, e.g., Serfling~\citet[Section 5.3]{serfling2009approximation}) gives
\[
\frac{1}{n(n-1)}\sum_{i \neq j} h_r(U_i, U_j)
= \theta + \frac{2}{n}\sum_{i=1}^n h_1(U_i) + U_n^{(2)},
\]
with 
\begin{itemize} \item $\theta := \mathbb{E}[h_r(U,U')] = \lambda_r$ (as computed in the proof of Lemma~\ref{lem:degU}); \item $h_1(x) := \mathbb{E}[h_r(x,U)] - \theta = \lambda_r\varphi_r(x)^2 - \lambda_r = \lambda_r(\varphi_r(x)^2 - 1)$; \item $U_n^{(2)} := \frac{1}{n(n-1)}\sum_{i \neq j} h_{r,2}(U_i, U_j)$ is the degenerate (second-order) U-statistic with $\mathbb{E}[U_n^{(2)}] = 0$ and $\Var(U_n^{(2)}) = O(n^{-2})$ (by standard U-statistic variance formulas for bounded kernels). \end{itemize}  Thus   \[ \frac{1}{n(n-1)}\sum_{i \neq j} h_r(U_i, U_j) = \lambda_r + \frac{2\lambda_r}{n}\sum_{i=1}^n (\varphi_r(U_i)^2 - 1) + U_n^{(2)}. \]   Define   \[ V_{r,n} := \frac{1}{n}\sum_{i=1}^n (\varphi_r(U_i)^2 - 1), \]   so that   \[ \frac{1}{n(n-1)}\sum_{i \neq j} h_r(U_i, U_j) = \lambda_r + 2\lambda_r V_{r,n} + U_n^{(2)}. \]
By Lemma~\ref{lem:norm}, $s_{r,n} = n(1 + V_{r,n})$, hence
\[
\frac{1}{s_{r,n}} = \frac{1}{n}(1 - V_{r,n} + O_p(n^{-1})).
\]
Therefore,
\begin{align*}
\mathbf{u}_r^\top \mathbf{K}_n \mathbf{u}_r
&= \frac{n(n-1)}{s_{r,n}} \cdot \frac{1}{n(n-1)}\sum_{i \neq j} h_r(U_i, U_j) \\
&= (n-1)(1 - V_{r,n} + O_p(n^{-1})) \cdot (\lambda_r + 2\lambda_r V_{r,n} + U_n^{(2)}) \\
&= (n-1)\lambda_r + (n-1)\lambda_r V_{r,n} + (n-1)U_n^{(2)} + O_p(1),
\end{align*}
where the cross-term $2\lambda_r V_{r,n} - \lambda_r V_{r,n} = \lambda_r V_{r,n}$ yields the
variance cancellation.

Now by construction, $V_{r,n}$ is the average of i.i.d.\ centered variables
$\varphi_r(U_i)^2 - 1$ with
\[
\mathbb{E}[\varphi_r(U)^2 - 1] = 0,\qquad
\Var(\varphi_r(U)^2 - 1) = \Var(\varphi_r(U)^2) = \sigma_r^2 > 0.
\]
Hence, by the classical central limit theorem,
\[
\sqrt{n} V_{r,n} \xrightarrow{d} \mathcal{N}(0, \sigma_r^2).
\]
For the degenerate U-statistic term, $\Var(U_n^{(2)}) = O(n^{-2})$ implies
\[
(n-1)U_n^{(2)} = O_p(1) = o_p(\sqrt{n}).
\]
Combining this with the above expansion of
$\mathbf{u}_r^\top \mathbf{K}_n \mathbf{u}_r$ and the approximation
$\lambda_r(\mathbf{K}_n) = \mathbf{u}_r^\top \mathbf{K}_n \mathbf{u}_r + o_p(\sqrt n)$, we obtain
\[
\lambda_r(\mathbf{K}_n) = (n-1)\lambda_r + (n-1)\lambda_r V_{r,n} + o_p(\sqrt{n}).
\]
Dividing by $\sqrt{n}$,
\[
\frac{\lambda_r(\mathbf{K}_n) - (n-1)\lambda_r}{\sqrt{n}}
= \lambda_r \sqrt{n} V_{r,n} + o_p(1)
\xrightarrow{d} \mathcal{N}(0, \lambda_r^2 \sigma_r^2),
\]
or equivalently,
\[
\sqrt{n}\Bigl(\frac{\lambda_r(\mathbf{K}_n)}{n-1} - \lambda_r\Bigr)
\xrightarrow{d} \mathcal{N}(0, \lambda_r^2 \sigma_r^2),
\]
as claimed.



\textbf{Case 2: Degenerate case ($\sigma_r^2 = 0$, equivalently, $\varphi_r^2 \equiv 1$ a.e.)}
In the degenerate case $\sigma_r^2 = \Var(\varphi_r(U)^2) = 0$, the condition holds if and only if
$\varphi_r(U)^2 = \E[\varphi_r(U)^2] = 1$ almost surely, which is equivalent to
$\varphi_r^2 \equiv 1$ almost everywhere on $[0,1]$. Since $\varphi_r^2 \equiv 1$, we have
$s_{r,n} = \sum_{i=1}^n \varphi_r(U_i)^2 = n$ exactly, and we define the unit vector
\[
\mathbf{u} := n^{-1/2}(\varphi_r(U_1),\dots,\varphi_r(U_n))^\top,
\qquad
\|\mathbf{u}\|^2 = 1.
\]
Let $\mathbf{V} \in \mathbb{R}^{n \times (n-1)}$ have orthonormal columns spanning $\mathbf{u}^\perp$.
By standard perturbation theory for simple eigenvalues (or Lemma~\ref{lem:expansion}), the eigenvalue
$\lambda_r(\mathbf{K}_n)$ then admits the expansion
\begin{equation}\label{expansion}
\begin{split}
\lambda_r(\mathbf{K}_n) &= 
\mathbf{u}^\top \mathbf{K}_n \mathbf{u} \\
&\quad + \mathbf{u}^\top \mathbf{K}_n \mathbf{V} 
\Bigl(
(n-1)\lambda_r \mathbf{I}_{n-1} - \mathbf{V}^\top \mathbf{K}_n \mathbf{V}
\Bigr)^{-1} 
\mathbf{V}^\top \mathbf{K}_n \mathbf{u}.
\end{split}
\end{equation}
For the first term in 
(\ref{expansion}), using the kernel representation and the definition
$h_r(x,y) = \varphi_r(x)W(x,y)\varphi_r(y)$,
\[
\mathbf{u}^\top \mathbf{K}_n \mathbf{u}
= \frac{1}{n}\sum_{i \neq j}\varphi_r(U_i) W(U_i, U_j) \varphi_r(U_j)
= \frac{1}{n}\sum_{i \neq j}h_r(U_i, U_j),
\]
and the Hoeffding decomposition $h_r = h_{r,1} + h_{r,2}$ with $h_{r,1} \equiv \lambda_r$
(Proof of Lemma~\ref{lem:degU}) yields
\[
\frac{1}{n}\sum_{i \neq j}h_r(U_i, U_j)
= \frac{n-1}{n}\lambda_r + \frac{1}{n}\sum_{i \neq j}h_{r,2}(U_i, U_j).
\]
By Lemma~\ref{lem:degU},
\[
\frac{1}{n}\sum_{i \neq j}h_{r,2}(U_i, U_j)
= \sum_{k \neq r}\lambda_k(T_{k,n}^2 - 1) + o_p(1),
\]
where $T_{k,n} := n^{-1/2}\sum_i \varphi_r(U_i)\varphi_k(U_i)$ and $(T_{k,n})$ converges jointly to
i.i.d.\ $\mathcal{N}(0,1)$ variables. Consequently,
\[
\mathbf{u}^\top \mathbf{K}_n \mathbf{u}
= (n-1)\lambda_r + \sum_{k \neq r}\lambda_k(T_{k,n}^2 - 1) + o_p(1).
\]
For the second term in (\ref{expansion}), Lemma~\ref{lem:resolvent} gives
\[
\mathbf{u}^\top \mathbf{K}_n \mathbf{V}\bigl((n-1)\lambda_r \mathbf{I}_{n-1}
- \mathbf{V}^\top\mathbf{K}_n\mathbf{V}\bigr)^{-1}\mathbf{V}^\top\mathbf{K}_n \mathbf{u}
= \sum_{k \neq r}\frac{\lambda_k^2}{\lambda_r - \lambda_k}T_{k,n}^2 + o_p(1).
\]
Combining these two contributions in the expression 
(\ref{expansion}), one obtains
\begin{align*}
\lambda_r(\mathbf{K}_n) &- (n-1)\lambda_r 
= \sum_{k \neq r}\lambda_k(T_{k,n}^2 - 1)
+ \sum_{k \neq r}\frac{\lambda_k^2}{\lambda_r - \lambda_k}T_{k,n}^2 + o_p(1) \\
&= \sum_{k \neq r}\lambda_k(T_{k,n}^2 - 1)
   + \sum_{k\neq r}\frac{\lambda_k^2}{\lambda_r-\lambda_k}(T_{k,n}^2-1)
   + \sum_{k\neq r}\frac{\lambda_k^2}{\lambda_r-\lambda_k}
   + o_p(1) \\
&= \sum_{k \neq r}\left[\lambda_k + \frac{\lambda_k^2}{\lambda_r - \lambda_k}\right](T_{k,n}^2-1)
   + \sum_{k\neq r}\frac{\lambda_k^2}{\lambda_r-\lambda_k}
   + o_p(1).
\end{align*}
Simplifying the coefficient,
\[
\lambda_k + \frac{\lambda_k^2}{\lambda_r - \lambda_k}
= \frac{\lambda_k(\lambda_r - \lambda_k) + \lambda_k^2}{\lambda_r - \lambda_k}
= \frac{\lambda_k\lambda_r}{\lambda_r - \lambda_k},
\]
we arrive at
\[
\lambda_r(\mathbf{K}_n) - (n-1)\lambda_r
=
\sum_{k \neq r}\frac{\lambda_r\lambda_k}{\lambda_r - \lambda_k}(T_{k,n}^2 - 1)
+ C_r + o_p(1),
\]
where the deterministic constant
\[
C_r := \sum_{k\neq r}\frac{\lambda_k^2}{\lambda_r-\lambda_k}
\]
is finite by the Hilbert-Schmidt condition and the eigengap. Thus the centered fluctuation can be written as
\[
\lambda_r(\mathbf{K}_n) - (n-1)\lambda_r - C_r
=
\sum_{k \neq r}\frac{\lambda_r\lambda_k}{\lambda_r - \lambda_k}(T_{k,n}^2 - 1)
+ o_p(1).
\]
Using the joint convergence from the proof of Lemma~\ref{lem:degU}, for any fixed $K$ the vector
$(T_{k,n})_{k \le K, k \neq r}$ converges jointly in distribution to $(Z_k)_{k \le K, k \neq r}$,
where $(Z_k)$ are i.i.d.\ $\mathcal{N}(0,1)$. By the Hilbert-Schmidt condition
$\sum_{k \ge 1}\lambda_k^2 < \infty$ and the eigengap bound $|\lambda_r - \lambda_k| \ge \gamma_r$
for all $k \neq r$,
\[
\sum_{k \neq r}\left(\frac{\lambda_r\lambda_k}{\lambda_r - \lambda_k}\right)^2\Var(Z_k^2-1)
=
2\sum_{k \neq r}\left(\frac{\lambda_r\lambda_k}{\lambda_r - \lambda_k}\right)^2
\le
\frac{2\lambda_r^2}{\gamma_r^2}\sum_{k\neq r}\lambda_k^2
<\infty,
\]
so the series
\[
\sum_{k \neq r}\frac{\lambda_r\lambda_k}{\lambda_r - \lambda_k}(Z_k^2 - 1)
\]
converges in $L^2$. Together with the $L^2$ tail control in Lemmas~\ref{lem:degU} and~\ref{lem:resolvent}, this implies
\[
\lambda_r(\mathbf{K}_n) - (n-1)\lambda_r - C_r
\xrightarrow{d}
\sum_{k \neq r}\frac{\lambda_r\lambda_k}{\lambda_r - \lambda_k}(Z_k^2 - 1),
\]
with convergence of the limit series in $L^2$, and, upon absorbing the deterministic second-order bias $C_r$ into the centering, one obtains the stated non-Gaussian limit in the degenerate case.


\end{proof}


\section*{Acknowledgments}
The author thanks Mohammad S.\,M.\ Moakhar for insightful discussions and valuable suggestions.


% =========================
% Appendix A: Spectral Perturbation Lemmas
% =========================

\appendix


\section*{Appendix A: Spectral Perturbation Lemmas}
\renewcommand{\thesection}{A.\arabic{section}}
\renewcommand{\thelemma}{A.\arabic{lemma}}

Throughout this appendix, Assumption~\ref{ass:ustat} is in force. Recall that
$W\in L^2([0,1]^2)\cap L^\infty([0,1]^2)$, that $T_W$ is the associated integral
operator, and that $(\lambda_k,\varphi_k)$ denote its eigenpairs.

\begin{lemma}[Approximate eigenpair for the target index $r$]
\label{lem:approx_eigvec}
Under Assumption~\ref{ass:ustat}, define
\[
\mathbf{u}_r := s_{r,n}^{-1/2}\bigl(\varphi_r(U_1),\dots,\varphi_r(U_n)\bigr)^\top,
\qquad
s_{r,n} := \sum_{i=1}^n \varphi_r(U_i)^2 .
\]
Then
\[
\|\mathbf{K}_n \mathbf{u}_r - (n-1)\lambda_r \mathbf{u}_r\|_{2} = O_p(\sqrt{n}),
\qquad
s_{r,n} = n(1 + o_p(1)).
\]
\end{lemma}


\begin{proof}
First, since $\E[\varphi_r(U)^2]=1$ and $\{\varphi_r(U_i)^2\}_{i=1}^n$ are i.i.d.,
the law of large numbers gives
\[
s_{r,n}=\sum_{i=1}^n \varphi_r(U_i)^2 = n(1+o_p(1)).
\]

Define the unnormalized vector
\[
\tilde{\mathbf u}_r := n^{-1/2}\bigl(\varphi_r(U_1),\dots,\varphi_r(U_n)\bigr)^\top .
\]
For each $i$,
\[
\Bigl(\frac{\mathbf K_n}{n-1}\tilde{\mathbf u}_r\Bigr)_i
=
\frac{1}{\sqrt n}\cdot \frac{1}{n-1}\sum_{j\ne i} W(U_i,U_j)\varphi_r(U_j).
\]
Let
\[
\Delta_i
:=
\frac{1}{n-1}\sum_{j\ne i} W(U_i,U_j)\varphi_r(U_j)-\lambda_r\varphi_r(U_i).
\]
Then $\E[\Delta_i\mid U_i]=0$. Conditioning on $U_i$, the summands are independent in $j$ and
\[
\E\!\left[ W(U_i,U_j)\varphi_r(U_j)\mid U_i\right]
=(T_W\varphi_r)(U_i)=\lambda_r\varphi_r(U_i).
\]
Moreover,
\[
\Var\!\left(W(U_i,U_j)\varphi_r(U_j)\mid U_i\right)
\le
\E\!\left[W(U_i,U_j)^2\varphi_r(U_j)^2\mid U_i\right]
\le
\|W\|_\infty^2\,\E[\varphi_r(U)^2]
=
\|W\|_\infty^2.
\]
Hence, conditional on $U_i$,
\[
\E[\Delta_i^2\mid U_i]
=
\Var(\Delta_i\mid U_i)
=
\frac{1}{(n-1)^2}\sum_{j\ne i}\Var\!\left(W(U_i,U_j)\varphi_r(U_j)\mid U_i\right)
\le \frac{C}{n-1},
\]
for a constant $C>0$. Therefore,
\[
\E\left\|\frac{\mathbf K_n}{n-1}\tilde{\mathbf u}_r-\lambda_r\tilde{\mathbf u}_r\right\|_2^2
=
\sum_{i=1}^n \E\left[\left(\frac{1}{\sqrt n}\Delta_i\right)^2\right]
=
\frac{1}{n}\sum_{i=1}^n \E[\Delta_i^2]
\le \frac{C}{n}.
\]
Thus
\[
\left\|\frac{\mathbf K_n}{n-1}\tilde{\mathbf u}_r-\lambda_r\tilde{\mathbf u}_r\right\|_2
=O_p(n^{-1/2}),
\]
equivalently,
\[
\|\mathbf K_n\tilde{\mathbf u}_r-(n-1)\lambda_r\tilde{\mathbf u}_r\|_2
=O_p(\sqrt n).
\]
Finally, since $\mathbf u_r=\sqrt{n/s_{r,n}}\,\tilde{\mathbf u}_r=(1+o_p(1))\tilde{\mathbf u}_r$,
the same bound holds with $\mathbf u_r$ in place of $\tilde{\mathbf u}_r$.
\end{proof}


\begin{lemma}[Sample eigengap bound]\label{lem:eigengap}
Under Assumption~\ref{ass:ustat}, let $T_W$ denote the integral operator with
eigenvalues $(\lambda_k)_{k\ge1}$, and assume that $\lambda_r$ is simple and isolated:
\[
\gamma_r := \min_{k\neq r}|\lambda_r-\lambda_k|>0.
\]
Let $\mathbf{K}_n$ be the kernel matrix and set
\[
\mathbf{K}_n^\star := \E\!\left[\mathbf{K}_n\,\middle|\,U_1,\dots,U_n\right].
\]
Assume further that
\[
\max_{1\le k\le n}\left|\frac{\lambda_k(\mathbf K_n^\star)}{n-1}-\lambda_k\right|=o_p(1).
\]
Then there exists a constant $c>0$ such that
\[
\mathbb{P}\!\left(
\min_{k\neq r}
\big|\lambda_r(\mathbf{K}_n) - \lambda_k(\mathbf{K}_n)\big|
\ge c n
\right)\to 1.
\]
\end{lemma}

\begin{proof}
By Lemma~\ref{lem:matconc},
\[
\Big\|\frac{1}{n-1}\mathbf{K}_n - \frac{1}{n-1}\mathbf{K}_n^\star\Big\|_{\op} = O_p(n^{-1/2}).
\]
By Weyl's inequality for self-adjoint matrices, for each $k\in[n]$,
\[
\Big|
\frac{\lambda_k(\mathbf{K}_n)}{n-1} - \frac{\lambda_k(\mathbf{K}_n^\star)}{n-1}
\Big|
\le
\Big\|\frac{1}{n-1}\mathbf{K}_n - \frac{1}{n-1}\mathbf{K}_n^\star\Big\|_{\op}.
\]
Combining with the assumption
$\max_{1\le k\le n}\left|\lambda_k(\mathbf K_n^\star)/(n-1)-\lambda_k\right|=o_p(1)$,
we obtain that with probability $1-o(1)$,
\[
\Big|
\frac{\lambda_k(\mathbf{K}_n)}{n-1} - \lambda_k
\Big|
\le \frac{\gamma_r}{4}
\qquad\text{for all }k\in[n].
\]
Therefore, on this event,
\[
\min_{k\neq r}
\Big|
\frac{\lambda_r(\mathbf{K}_n)}{n-1} - \frac{\lambda_k(\mathbf{K}_n)}{n-1}
\Big|
\ge
\min_{k\neq r}|\lambda_r-\lambda_k| - 2\cdot\frac{\gamma_r}{4}
=
\frac{\gamma_r}{2}.
\]
Multiplying by $(n-1)\ge n/2$ for large $n$ yields
\[
\min_{k\neq r}
|\lambda_r(\mathbf{K}_n) - \lambda_k(\mathbf{K}_n)|
\ge
\frac{\gamma_r}{4}\,n
\]
with probability $1-o(1)$, completing the proof.
\end{proof}







\begin{lemma}[Kato-Temple inequality for symmetric matrices]
\label{lem:kato-temple}
Let $\mathbf{A}\in\mathbb{R}^{n\times n}$ be symmetric, and let
$\mathbf{u}\in\mathbb{R}^n$ satisfy $\|\mathbf{u}\|_2=1$.
Define the Rayleigh quotient and residual
\[
\eta := \mathbf{u}^\top \mathbf{A}\mathbf{u},
\qquad
\mathbf{r} := \mathbf{A}\mathbf{u}-\eta\,\mathbf{u}.
\]
Suppose there exist real numbers $\alpha<\eta<\beta$ such that the open interval
$(\alpha,\beta)$ contains exactly one eigenvalue $\lambda$ of $\mathbf{A}$
(counting multiplicity).
Then
\[
\eta - \frac{\|\mathbf{r}\|_2^2}{\eta-\alpha}
\le
\lambda
\le
\eta + \frac{\|\mathbf{r}\|_2^2}{\beta-\eta}.
\]
In particular, if $\eta-\alpha$ and $\beta-\eta$ are of order $n$ and
$\|\mathbf{r}\|_2=o(\sqrt{n})$, then $|\lambda-\eta|=o(1)$.
\end{lemma}

\begin{proof}
This is a finite-dimensional specialization of the Kato--Temple inequality.
Specifically, equation (10) in \citet{Kato1949} states that if the interval
$(\alpha,\beta)$ contains no point of the spectrum except at most one
(non-degenerate) eigenvalue and
\[
\varepsilon^2 < (\eta-\alpha)(\beta-\eta),
\]
then
\[
\eta - \frac{\varepsilon^2}{\beta-\eta}
\le
\lambda
\le
\eta + \frac{\varepsilon^2}{\eta-\alpha}.
\]
Identifying $\varepsilon=\|\mathbf r\|_2$ yields the stated bounds.
\end{proof}



\begin{lemma}[Kernel matrix eigenvalue scaling in $(n-1)$ normalization]
\label{lem:matconc}
Under Assumption~\ref{ass:ustat}, let
\[
(\mathbf{K}_n)_{ij} := W(U_i,U_j)\,\mathbbm{1}\{i\neq j\},
\qquad
U_1,\dots,U_n \stackrel{\mathrm{iid}}{\sim}\mathrm{Unif}[0,1].
\]
Let $T_W$ be the integral operator on $L^2([0,1])$ with kernel $W$ and eigenvalues
$\lambda_1(T_W)\ge \lambda_2(T_W)\ge \cdots$.
Fix $r$ as in Assumption~\ref{ass:ustat}. Then
\[
\lambda_r(\mathbf{K}_n) = (n-1)\lambda_r + O_p(\sqrt n).
\]
In particular,
\[
\frac{\lambda_r(\mathbf{K}_n)}{n-1}\xrightarrow{p}\lambda_r,
\qquad
\text{and}
\qquad
\left|\frac{\lambda_r(\mathbf{K}_n)}{n-1}-\lambda_r \right|=O_p(n^{-1/2}).
\]
\end{lemma}

\begin{proof}
Let $\widetilde{\mathbf K}_n$ be the full kernel matrix with entries
$(\widetilde{\mathbf K}_n)_{ij}=W(U_i,U_j)$.
Then
\[
\mathbf K_n
=
\widetilde{\mathbf K}_n
-
\mathrm{diag}\big(W(U_1,U_1),\dots,W(U_n,U_n)\big).
\]
Since the difference is diagonal,
\[
\|\mathbf K_n-\widetilde{\mathbf K}_n\|_{\op}
=
\max_{1\le i\le n}|W(U_i,U_i)|
\le \|W\|_\infty,
\]
and therefore
\[
\Big\|\frac{1}{n-1}\mathbf K_n-\frac{1}{n-1}\widetilde{\mathbf K}_n\Big\|_{\op}
\le \frac{\|W\|_\infty}{n-1}
= o(n^{-1/2}).
\]

We view $(1/n)\widetilde{\mathbf K}_n$ as the empirical kernel integral operator associated
with the empirical measure $n^{-1}\sum_{i=1}^n \delta_{U_i}$.
By Koltchinskii and Giné (2000, Theorem~2.1), for a bounded measurable kernel
$W\in L^\infty([0,1]^2)$ and i.i.d.\ design points $U_i\sim\mathrm{Unif}[0,1]$,
\[
\Big\|\frac{1}{n}\widetilde{\mathbf K}_n-T_W\Big\|_{\op}
=O_p(n^{-1/2}).
\]

We next compare the normalizations $1/n$ and $1/(n-1)$:
\[
\Big\|\frac{1}{n-1}\widetilde{\mathbf K}_n-\frac{1}{n}\widetilde{\mathbf K}_n\Big\|_{\op}
=
\Big|\frac{1}{n-1}-\frac{1}{n}\Big|\,\|\widetilde{\mathbf K}_n\|_{\op}.
\]
Since $|(\widetilde{\mathbf K}_n)_{ij}|\le \|W\|_\infty$,
$\|\widetilde{\mathbf K}_n\|_{\op}\le n\|W\|_\infty$, and thus
\[
\Big\|\frac{1}{n-1}\widetilde{\mathbf K}_n-\frac{1}{n}\widetilde{\mathbf K}_n\Big\|_{\op}
\le \frac{\|W\|_\infty}{n-1}
=O(n^{-1})
=o(n^{-1/2}).
\]
Combining the last two displays yields
\[
\Big\|\frac{1}{n-1}\widetilde{\mathbf K}_n-T_W\Big\|_{\op}
=O_p(n^{-1/2}).
\]
Together with the diagonal-removal bound, we conclude
\[
\Big\|\frac{1}{n-1}\mathbf K_n-T_W\Big\|_{\op}
=O_p(n^{-1/2}).
\]

Finally, Weyl’s inequality (equivalently, the min-max characterization of eigenvalues)
applies to self-adjoint operators, giving for each fixed $r$,
\[
\Big|\frac{\lambda_r(\mathbf K_n)}{n-1}-\lambda_r\Big|
\le
\Big\|\frac{1}{n-1}\mathbf K_n-T_W\Big\|_{\op}
=O_p(n^{-1/2}).
\]
Multiplying both sides by $(n-1)$ gives
\[
\lambda_r(\mathbf K_n)
=
(n-1)\lambda_r
+
O_p(\sqrt n),
\]
which completes the proof.
\end{proof}



\begin{lemma}[Resolvent correction in the degenerate regime]\label{lem:resolvent}
Assume $\varphi_r^2 \equiv 1$ almost everywhere. Let 
\[
\mathbf{u} := n^{-1/2}(\varphi_r(U_1),\dots,\varphi_r(U_n))^\top
\]
and let $\mathbf{V} \in \mathbb{R}^{n \times (n-1)}$ have orthonormal 
columns spanning $\mathbf{u}^\perp$. 
Let $T_{k,n} := n^{-1/2}\sum_i \varphi_r(U_i)\varphi_k(U_i)$ for $k \neq r$. Then
\[
\mathbf{u}^\top \mathbf{K}_n \mathbf{V}\bigl((n-1)\lambda_r \mathbf{I}_{n-1} - \mathbf{V}^\top\mathbf{K}_n\mathbf{V}\bigr)^{-1}\mathbf{V}^\top\mathbf{K}_n \mathbf{u}
=
\sum_{k \neq r}\frac{\lambda_k^2}{\lambda_r - \lambda_k}T_{k,n}^2 + o_p(1).
\]
\end{lemma}

\begin{proof}
We approximate the sandwich product using the discretized eigenvectors $\{\mathbf u_k\}$ and a finite-rank truncation, then control the tail via $\sum_k\lambda_k^2<\infty$.

Fix the approximate eigenvectors $\mathbf u_k$ as in Lemma~\ref{lem:approx_eigvec}, and for each fixed $K$ consider the family $\{\mathbf u_k\}_{k\le K,k\neq r}$, which is asymptotically orthonormal and lies in $\mathbf u^\perp$ up to $o_p(1)$ by Lemma~\ref{lem:approx_eigvec}. By Gram-Schmidt, choose $\mathbf V$ so that, for each fixed $K$, its first $K-1$ columns agree with an orthonormalization of $\{\mathbf u_k\}_{k\le K,k\neq r}$, noting that the quadratic form
\[
\mathbf{u}^\top \mathbf{K}_n \mathbf{V}A^{-1}\mathbf{V}^\top\mathbf{K}_n \mathbf{u}
\]
depends only on the subspace $\mathbf u^\perp$ and $A$, not on the specific orthonormal basis $\mathbf V$, so this choice is without loss of generality. By Lemma~\ref{lem:matconc} and the simple eigengap at $\lambda_r$, the eigenvalues of $(n-1)^{-1}\mathbf K_n$ in $\mathbf u^\perp$ concentrate around $\{\lambda_k\}_{k\neq r}$, and hence
\[
\text{dist}\!\left(\lambda_r,\text{spec}\!\left(\frac{1}{n-1}\mathbf V^\top\mathbf K_n\mathbf V\right)\right)\ge \frac{\gamma_r}{2}
\]
with probability $1-o(1)$, so on this event
\[
\left\|\bigl((n-1)\lambda_r \mathbf{I}_{n-1} - \mathbf{V}^\top\mathbf{K}_n\mathbf{V}\bigr)^{-1}\right\|_{\op} 
\le \frac{2}{(n-1)\gamma_r}
=O(n^{-1}).
\]

For $k\ge1$, define
\[
\mathbf{u}_k:=s_{k,n}^{-1/2}\big(\varphi_k(U_1),\dots,\varphi_k(U_n)\big)^\top,
\qquad
s_{k,n}:=\sum_{i=1}^n\varphi_k(U_i)^2,
\]
so that $s_{k,n}=n(1+o_p(1))$ by Lemma~\ref{lem:norm}, and Lemma~\ref{lem:approx_eigvec} gives $\mathbf u_k^\top\mathbf u_\ell=\delta_{k\ell}+o_p(1)$ and
\[
\mathbf K_n\mathbf u_k=(n-1)\lambda_k\mathbf u_k+\mathbf e_k,
\qquad
\|\mathbf e_k\|=O_p(\sqrt n)
\quad\text{for each fixed }k.
\]
For $k\neq r$, the cross–projection scaling
\[
\mathbf u^\top\mathbf u_k
=
\frac{1}{\sqrt{ns_{k,n}}}\sum_{i=1}^n\varphi_r(U_i)\varphi_k(U_i)
=
\frac{1}{\sqrt n}\,(1+o_p(1))\,T_{k,n}
\]
holds. Fix $K<\infty$ with $K>r$ and define the finite–rank approximation
\[
\mathbf M_K:=\sum_{\substack{k\le K\\ k\neq r}}(n-1)\lambda_k\,\mathbf u_k\mathbf u_k^\top = \mathbf M_K=(n-1)\mathbf U\, \mathrm{diag}(\lambda_k)_{k\le K,\,k\neq r}\, \mathbf U^\top,
\]
where  $\mathbf U:=[\mathbf u_k]_{k\le K,\,k\neq r}\in\mathbb R^{n\times (K-1)}$.

In the approximately orthonormal basis $\{\mathbf u_k\}_{k\le K,\,k\neq r}$ we have

\[
\Big\|
\bigl((n-1)\lambda_r I-\mathbf M_K\bigr)^{-1}
-
\sum_{\substack{k\le K\\ k\neq r}}
\frac{1}{(n-1)(\lambda_r-\lambda_k)}\,\mathbf u_k\mathbf u_k^\top
\Big\|_{\op}
=o_p(n^{-1}),
\]
where we used lemma ~\ref{lem:approx_eigvec} for $\mathbf U^\top\mathbf U$ and $\mathbf{M}_K$.
For $k\le K$, $k\neq r$, the approximate eigenvalue equation yields
\[
\mathbf u^\top\mathbf K_n\mathbf u_k
=
(n-1)\lambda_k(\mathbf u^\top\mathbf u_k)+\mathbf u^\top\mathbf e_k
=
\frac{(n-1)\lambda_k}{\sqrt n}\,T_{k,n}\,(1+o_p(1)),
\]
since $|\mathbf u^\top\mathbf e_k|\le \|\mathbf e_k\|=O_p(\sqrt n)$. Therefore, sandwiching the truncated resolvent gives
\[
\mathbf u^\top\mathbf K_n\mathbf V\bigl((n-1)\lambda_r\mathbf I-\mathbf M_K\bigr)^{-1}\mathbf V^\top\mathbf K_n\mathbf u
=
\sum_{\substack{k\le K\\ k\neq r}}
\frac{\lambda_k^2}{\lambda_r-\lambda_k}\,T_{k,n}^2
+ o_p(1)
\]
To control the tail, use $|\lambda_r-\lambda_k|\ge \gamma_r$ for $k\neq r$ to get
\[
\left|\sum_{k>K}\frac{\lambda_k^2}{\lambda_r-\lambda_k}T_{k,n}^2\right|
\le \frac{1}{\gamma_r}\sum_{k>K}\lambda_k^2T_{k,n}^2.
\]
Since $\mathbb{E}[T_{k,n}^2]=1$ for each $k\neq r$, one has
\[
\mathbb{E}\!\left[\sum_{k>K}\lambda_k^2T_{k,n}^2\right]
=
\sum_{k>K}\lambda_k^2
\xrightarrow[K\to\infty]{}0,
\]
and Markov's inequality for nonnegative random variables then implies
\[
\lim_{K\to\infty}\sup_n\mathbb{P}\!\left(\left|\sum_{k>K}\frac{\lambda_k^2}{\lambda_r-\lambda_k}T_{k,n}^2\right|>\varepsilon\right)=0,
\]
so combining the finite–mode identity with this tail bound and letting $K\to\infty$ yields the result.
\end{proof}

\begin{lemma}[Second-order Rayleigh-Schrödinger expansion for a simple eigenvalue]
\label{lem:expansion}
Let $\mathbf{M} \in \mathbb{R}^{n \times n}$ be symmetric, and fix $r \in \{1,\dots,n\}$ such that 
$\lambda_r := \lambda_r(\mathbf{M})$ is simple with eigengap
\[
\gamma_r := \min_{k \neq r}|\lambda_r(\mathbf{M}) - \lambda_k(\mathbf{M})| > 0.
\]
Define the perturbed matrix $\widehat{\mathbf{M}} = \mathbf{M} + \mathbf{E}$, and let $\mathbf{u}_r$
be the unit eigenvector of $\mathbf{M}$ corresponding to $\lambda_r$.
If $\|\mathbf{E}\|_{\op} < \gamma_r/2$, then
\[
\widetilde{\lambda}_r - \lambda_r
=
\mathbf{u}_r^\top \mathbf{E}\mathbf{u}_r + R,
\]
where $\widetilde{\lambda}_r := \lambda_r(\widehat{\mathbf{M}})$ and the remainder satisfies
\[
|R| \le \frac{2\|\mathbf{E}\|_{\op}^2}{\gamma_r}.
\]
\end{lemma}

\begin{proof}
We follow the block decomposition with respect to the eigenspace of $\lambda_r$.
Let $\{\mathbf{u}_1,\dots,\mathbf{u}_n\}$ be an orthonormal eigenbasis of $\mathbf{M}$ with
$\mathbf{M}\mathbf{u}_k=\lambda_k(\mathbf{M})\mathbf{u}_k$ and $\mathbf{u}_r$ corresponding to $\lambda_r$.
Define
\[
\mathbf{Q}
:=
[\mathbf{u}_1,\dots,\mathbf{u}_{r-1},\mathbf{u}_{r+1},\dots,\mathbf{u}_n]
\in \mathbb{R}^{n\times(n-1)},
\]
so that $\mathbf{Q}$ spans the orthogonal complement of $\mathbf{u}_r$ and
\(
\mathbf{Q}^\top \mathbf{Q}=\mathbf{I}_{n-1},
\,
\mathbf{Q}\mathbf{Q}^\top=\mathbf{I}_n-\mathbf{u}_r\mathbf{u}_r^\top.
\)

Every $\mathbf{x}\in\mathbb{R}^n$ admits a unique decomposition
\[
\mathbf{x}=\mathbf{u}_r\xi+\mathbf{Q}\mathbf{z},
\qquad
\xi\in\mathbb{R},\ \mathbf{z}\in\mathbb{R}^{n-1}.
\]

Let $\widetilde{\mathbf{u}}_r$ be the unit eigenvector of $\widehat{\mathbf{M}}$
associated with $\widetilde{\lambda}_r$, and write
\[
\widetilde{\mathbf{u}}_r=\mathbf{u}_r\xi+\mathbf{Q}\mathbf{z},
\qquad
\xi^2+\|\mathbf{z}\|^2=1.
\]

The eigenvalue equation
\[
(\mathbf{M}+\mathbf{E})\widetilde{\mathbf{u}}_r
=
\widetilde{\lambda}_r \widetilde{\mathbf{u}}_r
\]
becomes
\[
(\mathbf{M}+\mathbf{E})(\mathbf{u}_r\xi+\mathbf{Q}\mathbf{z})
=
\widetilde{\lambda}_r(\mathbf{u}_r\xi+\mathbf{Q}\mathbf{z}).
\]

Projecting onto $\mathbf{u}_r$ and onto $\text{span}(\mathbf{Q})$ yields
\[
\begin{cases}
\lambda_r\xi + \mathbf{u}_r^\top \mathbf{E}(\mathbf{u}_r\xi+\mathbf{Q}\mathbf{z})
= \widetilde{\lambda}_r \xi, \\[0.3em]
\mathbf{Q}^\top \mathbf{M}\mathbf{Q}\mathbf{z}
+ \mathbf{Q}^\top \mathbf{E}(\mathbf{u}_r\xi+\mathbf{Q}\mathbf{z})
= \widetilde{\lambda}_r \mathbf{z}.
\end{cases}
\]

Let $\alpha:=\widetilde{\lambda}_r-\lambda_r$.
Rewrite the second equation as
\[
\bigl(\mathbf{L}_0 + \mathbf{\Delta}\bigr)\mathbf{z}
=
-\xi\,\mathbf{Q}^\top\mathbf{E}\mathbf{u}_r,
\]
where
\[
\mathbf{L}_0 := \mathbf{Q}^\top(\mathbf{M}-\lambda_r\mathbf{I})\mathbf{Q},
\qquad
\mathbf{\Delta} := \mathbf{Q}^\top(\mathbf{E}-\alpha\mathbf{I})\mathbf{Q}.
\]

By the eigengap assumption, $\|\mathbf{L}_0^{-1}\|_{\op}\le 1/\gamma_r$.
Under $\|\mathbf{E}\|_{\op}<\gamma_r/2$,
$\mathbf{L}_0+\mathbf{\Delta}$ is invertible and
\[
\|(\mathbf{L}_0+\mathbf{\Delta})^{-1}\|_{\op}
\le \frac{2}{\gamma_r}.
\]

Hence
\[
\mathbf{z}
=
-\xi\,(\mathbf{L}_0+\mathbf{\Delta})^{-1}
\mathbf{Q}^\top\mathbf{E}\mathbf{u}_r.
\]

From the first projected equation,
\[
\alpha\xi
=
\xi\,\mathbf{u}_r^\top\mathbf{E}\mathbf{u}_r
+
\mathbf{u}_r^\top\mathbf{E}\mathbf{Q}\mathbf{z}.
\]

If $\xi\neq0$, divide by $\xi$ and substitute $\mathbf{z}$:
\[
\alpha
=
\mathbf{u}_r^\top\mathbf{E}\mathbf{u}_r
-
\mathbf{u}_r^\top\mathbf{E}\mathbf{Q}
(\mathbf{L}_0+\mathbf{\Delta})^{-1}
\mathbf{Q}^\top\mathbf{E}\mathbf{u}_r.
\]

We now expand the inverse via a resolvent identity:
\[
(\mathbf{L}_0+\mathbf{\Delta})^{-1}
=
\mathbf{L}_0^{-1}
-
\mathbf{L}_0^{-1}\mathbf{\Delta}\mathbf{L}_0^{-1}
+
\mathbf{R}_3,
\]
where
\(
\|\mathbf{R}_3\|_{\op}
\le C \|\mathbf{E}\|_{\op}^2/\gamma_r^3.
\)

Substituting gives
\[
\alpha
=
\mathbf{u}_r^\top\mathbf{E}\mathbf{u}_r
-
\mathbf{u}_r^\top\mathbf{E}\mathbf{Q}
\mathbf{L}_0^{-1}
\mathbf{Q}^\top\mathbf{E}\mathbf{u}_r
+
R_3,
\]
with
\(
|R_3|
\le C\|\mathbf{E}\|_{\op}^3/\gamma_r^2.
\)

Since
\[
\mathbf{Q}\mathbf{L}_0^{-1}\mathbf{Q}^\top
=
\sum_{k\ne r}
\frac{1}{\lambda_k-\lambda_r}
\mathbf{u}_k\mathbf{u}_k^\top,
\]
we obtain the explicit second-order Rayleigh–Schrödinger term
\[
\mathbf{u}_r^\top\mathbf{E}\mathbf{Q}
\mathbf{L}_0^{-1}
\mathbf{Q}^\top\mathbf{E}\mathbf{u}_r
=
\sum_{k\ne r}
\frac{(\mathbf{u}_k^\top\mathbf{E}\mathbf{u}_r)^2}
{\lambda_r-\lambda_k}.
\]

Therefore,
\[
\widetilde{\lambda}_r-\lambda_r
=
\mathbf{u}_r^\top\mathbf{E}\mathbf{u}_r
+
\sum_{k\ne r}
\frac{(\mathbf{u}_k^\top\mathbf{E}\mathbf{u}_r)^2}
{\lambda_r-\lambda_k}
+
R_3,
\]
with
\[
|R_3|
\le
C\frac{\|\mathbf{E}\|_{\op}^3}{\gamma_r^2}.
\]

This yields the full second-order Rayleigh–Schrödinger expansion.
\end{proof}

\section*{Appendix B: U-Statistic and Kernel Decompositions} \setcounter{section}{2} \renewcommand{\thesection}{B.\arabic{section}} \renewcommand{\thelemma}{B.\arabic{lemma}}

\begin{lemma}[Degenerate U-statistic in the degenerate regime]\label{lem:degU}
Assume $\varphi_r^2 \equiv 1$ almost everywhere. Let $h_r(x,y) := \varphi_r(x)W(x,y)\varphi_r(y)$ and define the Hoeffding-projected degenerate kernel
\[
h_{r,2}(x,y) := h_r(x,y) - \mathbb{E}[h_r(x,U)] - \mathbb{E}[h_r(U,y)] + \mathbb{E}[h_r(U,U')],
\]
where $U, U' \stackrel{\text{iid}}{\sim} \text{Unif}[0,1]$. Define the cross-projections
\[
T_{k,n} := \frac{1}{\sqrt{n}}\sum_{i=1}^n \varphi_r(U_i)\varphi_k(U_i), \qquad k \neq r.
\]
Then
\[
\frac{1}{n}\sum_{i \neq j}h_{r,2}(U_i,U_j)
=
\sum_{k \neq r}\lambda_k\left(T_{k,n}^2 - 1\right) + o_p(1),
\]
where the series converges in $L^2$-norm under the Hilbert-Schmidt condition $\sum_{k \ge 1}\lambda_k^2 < \infty$.
\end{lemma}

\begin{proof}
We first obtain an exact spectral identity, then pass to the infinite sum by truncation and an $L^2$ tail bound. Since $\varphi_r(x)^2=1$ a.e., we have $\varphi_r\in\{-1,1\}$ a.e. Using
$W(x,y)=\sum_{k\ge1}\lambda_k\varphi_k(x)\varphi_k(y)$,
\[
h_r(x,y)=\varphi_r(x)W(x,y)\varphi_r(y)
=\sum_{k\ge1}\lambda_k\,[\varphi_r(x)\varphi_k(x)]\,[\varphi_r(y)\varphi_k(y)].
\]
By orthonormality, $\mathbb{E}[\varphi_r(U)\varphi_k(U)]=\langle \varphi_r,\varphi_k\rangle_{L^2}=\delta_{rk}$, hence
\[
\mathbb{E}[h_r(x,U)]
=
\varphi_r(x)\sum_{k\ge1}\lambda_k\varphi_k(x)\,\delta_{rk}
=
\lambda_r\varphi_r(x)^2=\lambda_r.
\]
By symmetry, $\mathbb{E}[h_r(U,y)]=\lambda_r$ and $\mathbb{E}[h_r(U,U')]=\lambda_r$, so
\[
h_{r,2}(x,y)=h_r(x,y)-\lambda_r
=\varphi_r(x)\varphi_r(y)\sum_{k\neq r}\lambda_k\,\varphi_k(x)\varphi_k(y).
\]
For $k\neq r$, set $\xi_{k,i}:=\varphi_r(U_i)\varphi_k(U_i)$. Then
\begin{align*}
\frac{1}{n}\sum_{i\neq j}h_{r,2}(U_i,U_j)
&=
\sum_{k\neq r}\lambda_k\cdot \frac{1}{n}\sum_{i\neq j}\xi_{k,i}\xi_{k,j}\\
&=
\sum_{k\neq r}\lambda_k\left[
\frac{1}{n}\Big(\sum_{i=1}^n\xi_{k,i}\Big)^2-\frac{1}{n}\sum_{i=1}^n\xi_{k,i}^2
\right].
\end{align*}
Since $T_{k,n}=n^{-1/2}\sum_{i=1}^n\xi_{k,i}$ and $\xi_{k,i}^2=\varphi_r(U_i)^2\varphi_k(U_i)^2=\varphi_k(U_i)^2$,
\[
\frac{1}{n}\sum_{i\neq j}h_{r,2}(U_i,U_j)
=
\sum_{k\neq r}\lambda_k\left(
T_{k,n}^2-\frac{1}{n}\sum_{i=1}^n\varphi_k(U_i)^2
\right).
\]
Because $\mathbb{E}[\varphi_k(U)^2]=1$, the strong law gives
$n^{-1}\sum_{i=1}^n\varphi_k(U_i)^2\to 1$ a.s. for each fixed $k$.
Therefore, for any fixed $K<\infty$,
\[
\sum_{\substack{k\le K\\ k\neq r}}\lambda_k\left(
T_{k,n}^2-\frac{1}{n}\sum_{i=1}^n\varphi_k(U_i)^2
\right)
=
\sum_{\substack{k\le K\\ k\neq r}}\lambda_k(T_{k,n}^2-1)+o_p(1),
\]
where the $o_p(1)$ comes from a finite sum of terms
$\lambda_k\left(1-n^{-1}\sum_i\varphi_k(U_i)^2\right)$.

Fix $K<\infty$. For each $k\neq r$, $T_{k,n}$ is a normalized sum of i.i.d.\ centered variables.
Indeed, $\mathbb{E}[\varphi_r(U)\varphi_k(U)]=0$ for $k\neq r$, and
\[
\Var(\varphi_r(U)\varphi_k(U))
=
\mathbb{E}[\varphi_r(U)^2\varphi_k(U)^2]
=
\mathbb{E}[\varphi_k(U)^2]
=
1,
\]
since $\varphi_r^2\equiv 1$ almost everywhere.
Therefore, for each fixed $k\neq r$, the classical central limit theorem gives
$T_{k,n}\xrightarrow{d}\mathcal{N}(0,1)$.
Moreover, for any fixed coefficients $(a_k)_{k\le K,\,k\neq r}$,
the linear combination $\sum_{k\le K,\,k\neq r} a_k T_{k,n}$ is again a normalized sum of i.i.d.\ variables
with finite variance, so the Cram\'er-Wold device yields joint convergence of
$(T_{k,n})_{k\le K,\,k\neq r}$ to i.i.d.\ $\mathcal{N}(0,1)$ limits.

To control the tail, define the tail kernel and associated U-statistic
\[
h_{r,2}^{(K)}(x,y):=\varphi_r(x)\varphi_r(y)\sum_{k>K}\lambda_k\varphi_k(x)\varphi_k(y),
\qquad
U_n^{(K)}:=\frac{1}{n}\sum_{i\neq j} h_{r,2}^{(K)}(U_i,U_j).
\]
By the same algebra as above applied to the truncated series (restricting the sum to $k>K$), we also have the identity
\[
U_n^{(K)}=\sum_{k>K}\lambda_k\left(T_{k,n}^2-\frac{1}{n}\sum_{i=1}^n\varphi_k(U_i)^2\right).
\]
Since $h_{r,2}^{(K)}$ is symmetric and degenerate (i.e., $\mathbb{E}[h_{r,2}^{(K)}(x,U)]=0$ for all $x$), the standard second-moment identity for degenerate U-statistics (see Hoeffding \citet{hoeffding1992class}) yields a bound of the form
\[
\sup_{n\ge 2}\mathbb{E}\big[(U_n^{(K)})^2\big]\ \lesssim\ \|h_{r,2}^{(K)}\|_{L^2([0,1]^2)}^2.
\]
Using orthonormality of $\{\varphi_k\}$,
\[
\|h_{r,2}^{(K)}\|_{L^2}^2
=
\int_0^1\!\!\int_0^1\Big(\sum_{k>K}\lambda_k\varphi_k(x)\varphi_k(y)\Big)^2dx\,dy
=
\sum_{k>K}\lambda_k^2,
\]
so
\[
\sup_{n\ge 2}\mathbb{E}\big[(U_n^{(K)})^2\big]\ \lesssim\ \sum_{k>K}\lambda_k^2
\ \xrightarrow[K\to\infty]{}\ 0.
\]
By Markov's inequality,
$\lim_{K\to\infty}\sup_n\mathbb{P}(|U_n^{(K)}|>\varepsilon)=0$ for every $\varepsilon>0$.

Combining the exact identity
\[
\frac{1}{n}\sum_{i\neq j}h_{r,2}(U_i,U_j)
=
\sum_{k\neq r}\lambda_k\left(
T_{k,n}^2-\frac{1}{n}\sum_{i=1}^n\varphi_k(U_i)^2
\right)
\]
with the finite-$K$ approximation from the strong law and the joint CLT, and then sending $K\to\infty$ using the $L^2$ tail bound for $U_n^{(K)}$, we conclude that
\[
\frac{1}{n}\sum_{i\neq j}h_{r,2}(U_i,U_j)
=
\sum_{k\neq r}\lambda_k(T_{k,n}^2-1)+o_p(1),
\]
and the series convergence holds in $L^2$ (hence in probability).
\end{proof}

\begin{lemma}[Square-summability of fluctuation coefficients]\label{lem:squaresum}
Let \(T_W:L^2([0,1])\to L^2([0,1])\) be the compact self-adjoint operator associated with the symmetric kernel \(W\in L^2([0,1]^2)\cap L^\infty([0,1]^2)\), and let \(\{(\lambda_k,\varphi_k)\}_{k\ge1}\) be its orthonormal eigenpairs. 
Fix a simple eigenvalue \(\lambda_r\) with eigengap 
\(\gamma_r = \inf_{k\neq r} |\lambda_r - \lambda_k| > 0\).
Then
\[
\sum_{k\neq r} \left( \frac{\lambda_r\lambda_k}{\lambda_r - \lambda_k} \right)^2 < \infty.
\]
\end{lemma}

\begin{proof}
Since \(T_W\) is Hilbert-Schmidt, we have 
\(\sum_{k=1}^\infty \lambda_k^2 = \|W\|_{L^2([0,1]^2)}^2 < \infty.\)
For \(k\neq r\), note that \(|\lambda_r - \lambda_k| \ge \gamma_r > 0\) by the eigengap assumption, so
\[
\left(\frac{\lambda_r\lambda_k}{\lambda_r - \lambda_k}\right)^2 
\le \frac{\lambda_r^2}{\gamma_r^2}\,\lambda_k^2.
\]
Summing over all \(k\neq r\) yields
\[
\sum_{k\neq r}\left(\frac{\lambda_r\lambda_k}{\lambda_r - \lambda_k}\right)^2
\le \frac{\lambda_r^2}{\gamma_r^2} \sum_{k\neq r}\lambda_k^2
\le \frac{\lambda_r^2}{\gamma_r^2}\|W\|_{L^2([0,1]^2)}^2 < \infty.
\]
Hence the coefficients are square-summable, and the random series
\(\sum_{k\neq r} \tfrac{\lambda_r\lambda_k}{\lambda_r - \lambda_k}(Z_k^2-1)\)
converges in \(L^2\) by the completeness of \(\ell^2\) and independence of the Gaussian sequence \((Z_k)\).
\end{proof}







\begin{lemma}\label{lem:norm}
Let
\[
V_{r,n}:=\frac{1}{n}\sum_{i=1}^n(\varphi_r(U_i)^2-1),
\qquad
s_n:=\sum_{i=1}^n \varphi_r(U_i)^2 = n(1+V_{r,n}).
\]
Assume $\Var(\varphi_r(U)^2)<\infty$. Then $V_{r,n}=O_p(n^{-1/2})$ and
\[
\frac{1}{s_n}
=
\frac{1}{n}\Bigl(1 - V_{r,n} + O_p(V_{r,n}^2)\Bigr),
\qquad
\frac{1}{\sqrt{s_n}}
=
\frac{1}{\sqrt{n}}\Bigl(1 - \tfrac{1}{2}V_{r,n} + O_p(V_{r,n}^2)\Bigr).
\]
In particular, since $V_{r,n}^2=O_p(n^{-1})$, both remainders are $O_p(n^{-1})$.
\end{lemma}

\begin{proof}
Assume $\Var(\varphi_r(U)^2)<\infty$ and recall
\[
V_{r,n}:=\frac{1}{n}\sum_{i=1}^n\bigl(\varphi_r(U_i)^2-1\bigr),
\qquad
s_n:=\sum_{i=1}^n \varphi_r(U_i)^2=n(1+V_{r,n}).
\]
Since $\E[\varphi_r(U)^2]=1$ and $\Var(\varphi_r(U)^2)<\infty$, we have
\[
\E[V_{r,n}]=0,
\qquad
\Var(V_{r,n})=\frac{1}{n}\Var(\varphi_r(U)^2),
\]
hence by Chebyshev's inequality $V_{r,n}=O_p(n^{-1/2})$.

For the first expansion, write
\[
\frac{1}{s_n}=\frac{1}{n}(1+V_{r,n})^{-1}.
\]
Using the Taylor expansion of $f(x)=(1+x)^{-1}$ at $0$,
\[
(1+x)^{-1}=1-x+x^2+R_1(x),
\]
where the remainder satisfies $|R_1(x)|\le C|x|^3$ for all $|x|\le 1/2$ and a
universal constant $C>0$. On the event $\{|V_{r,n}|\le 1/2\}$,
\[
(1+V_{r,n})^{-1}
=
1-V_{r,n}+V_{r,n}^2+R_1(V_{r,n}),
\qquad
|R_1(V_{r,n})|\le C|V_{r,n}|^3.
\]
Since $V_{r,n}=O_p(n^{-1/2})$, we have $V_{r,n}^2=O_p(n^{-1})$ and
$V_{r,n}^3=O_p(n^{-3/2})$, so
\[
(1+V_{r,n})^{-1}
=
1-V_{r,n}+O_p(V_{r,n}^2).
\]
Multiplying by $1/n$ yields
\[
\frac{1}{s_n}
=
\frac{1}{n}\Bigl(1-V_{r,n}+O_p(V_{r,n}^2)\Bigr).
\]

For the second expansion, similarly
\[
\frac{1}{\sqrt{s_n}}=\frac{1}{\sqrt n}(1+V_{r,n})^{-1/2}.
\]
Using the Taylor expansion of $g(x)=(1+x)^{-1/2}$ at $0$,
\[
(1+x)^{-1/2}=1-\tfrac12 x+R_2(x),
\]
where $|R_2(x)|\le C'x^2$ for all $|x|\le 1/2$ and a universal constant $C'>0$.
On $\{|V_{r,n}|\le 1/2\}$ we obtain
\[
(1+V_{r,n})^{-1/2}
=
1-\tfrac12 V_{r,n}+R_2(V_{r,n}),
\qquad
|R_2(V_{r,n})|\le C'V_{r,n}^2,
\]
hence
\[
(1+V_{r,n})^{-1/2}
=
1-\tfrac12 V_{r,n}+O_p(V_{r,n}^2).
\]
Multiplying by $1/\sqrt n$ gives
\[
\frac{1}{\sqrt{s_n}}
=
\frac{1}{\sqrt n}\Bigl(1-\tfrac12 V_{r,n}+O_p(V_{r,n}^2)\Bigr).
\]
Since $V_{r,n}^2=O_p(n^{-1})$, both $O_p(V_{r,n}^2)$ remainders are also
$O_p(n^{-1})$.
\end{proof}








\bibliographystyle{apalike}
\bibliography{Reference}
\end{document}
